\chapter{כלכלת הקשב}
\label{ch:attention}

\begin{epigraph}{הרברט סיימון, 1971}
מה שמידע צורך הוא די ברור: הוא צורך את תשומת הלב של מקבליו.
\end{epigraph}

\section{עוני של קשב}

במשך רוב ההיסטוריה, מידע היה משאב נדיר. ספר היה אוצר. מכתב היה אירוע. מי ששלט במידע, הכוהנים, הסופרים והשליטים, שלט בחברה. הבעיה של בני האדם הייתה תמיד מחסור במידע.

ב-1971, הרבה לפני הסמארטפון ואפילו לפני המחשב האישי, הבחין הכלכלן והמדען הרברט סיימון שהמשוואה עומדת להתהפך. מידע, כתב סיימון, איננו בחינם. הוא צורך משהו: את תשומת הלב של מי שקולט אותו. ומכאן המסקנה שלו, שנשמעת היום כמו נבואה: \textbf{שפע של מידע יוצר עוני של קשב} \parencite{simon1971}.

חשבו על זה רגע. היום יש בכיס שלכם יותר מידע מאשר בכל הספריות של העולם העתיק גם יחד. אבל יש לכם בדיוק אותן 24 שעות ביממה שהיו לאריסטו. המידע הפך לזול כמו מים. הקשב נשאר יקר כמו זהב. וכשמשהו נדיר ויקר, מישהו תמיד ימצא דרך לסחור בו.

\section{סוחרי הקשב}

הרעיון למכור קשב לא נולד בעמק הסיליקון. טים וו, משפטן והיסטוריון של התקשורת, מתאר בספרו "סוחרי הקשב" כיצד העיקרון התגבש כבר במאה ה-19, עם עיתוני הפרוטה הזולים, ועבר משם לרדיו, לטלוויזיה ולבסוף לרשתות החברתיות \parencite{wu2016}. המודל העסקי פשוט: נותנים לקהל תוכן בחינם או כמעט בחינם, אוספים את תשומת הלב שלו, ומוכרים אותה למפרסמים.

ההבדל בין עיתון הפרוטה לבין הפיד שלכם איננו בעיקרון אלא בדיוק. העיתון ידע בערך מי קורא אותו. האפליקציה יודעת בדיוק כמה שניות התעכבתם על כל תמונה.

ויש הבדל נוסף. עיתון נגמר. בעמוד האחרון יש תשבץ, ואחריו אין כלום. גם שידורי הטלוויזיה היו נגמרים פעם בחצות, עם המנון ופס צבעים. הפיד, לעומת זאת, איננו נגמר לעולם. אין לו עמוד אחרון, ולכן ההחלטה מתי להפסיק עוברת כולה אלינו, בדיוק בשעה שבה אנחנו הכי עייפים.

\begin{factbox}[המשפט שכדאי לזכור]
אם אינכם משלמים על המוצר, סביר להניח שאתם לא הלקוחות. אתם חומר הגלם. הלקוחות הם המפרסמים, והמוצר הוא הקשב שלכם.
\end{factbox}

\section{מודל ה-\texorpdfstring{\en{Hook}}{Hook}}

איך גורמים לאדם לחזור לאפליקציה עשרות פעמים ביום, בלי שאיש מכריח אותו? ניר אייל, יועץ ומחבר הספר \en{Hooked}, תיאר ב-2014 מודל לבניית מוצרים שיוצרים הרגל \parencite{eyal2014}. המודל בנוי מארבעה שלבים שחוזרים על עצמם במעגל: \textbf{טריגר}, \textbf{פעולה}, \textbf{תגמול משתנה} ו\textbf{השקעה}.

הלב של המודל הוא התגמול ה\textbf{משתנה}. אילו כל פתיחה של האפליקציה הייתה מניבה בדיוק אותו דבר, היינו משתעממים מהר. אבל כשאיננו יודעים מה מחכה לנו, לייק, הודעה מאדם שאנחנו אוהבים או סתם כלום, אנחנו חוזרים שוב ושוב. מכונות המזל בקזינו עובדות על אותו עיקרון בדיוק. ההבדל הוא שאת מכונת המזל הזאת אנחנו נושאים בכיס, והיא פתוחה עשרים וארבע שעות ביממה.

\section{מעגל בן שמונה שלבים}

מודל ה-\en{Hook} מסביר מה קורה בצד של המשתמש. אבל כדי להבין את כלכלת הקשב, צריך לראות גם את הצד השני של המסך. איור~\ref{fig:attention-loop} מחבר את שני הצדדים למעגל אחד בן שמונה שלבים.

% ---------- תרשים בלוקים: מעגל כלכלת הקשב (8 בלוקים) ----------
\begin{figure}[htbp]
\centering
\begin{tikzpicture}[
  blk/.style={rectangle, rounded corners=5pt, draw=#1, line width=1.2pt,
              fill=#1!12, text width=3.1cm, minimum height=1.25cm,
              align=center, font=\headfont\small, drop shadow={opacity=0.15}},
  blk/.default=screen,
  arr/.style={-{Stealth[length=3mm]}, line width=1.4pt, draw=night!70},
  num/.style={circle, fill=#1, text=white, font=\headfont\bfseries\footnotesize,
              inner sep=1.5pt, minimum size=1.3em},
]
  % מרכז
  \node[circle, fill=night, text=white, minimum size=2.6cm, align=center,
        font=\headfont\bfseries, text width=1.8cm] (core) at (0,0) {\RL{הקשב}\\ \RL{שלך}};

  % 8 בלוקים במעגל (נקראים עם כיוון השעון, מלמעלה)
  % \RL{} שומר על סדר המילים העברי בתוך צמתי TikZ
  \def\rad{4.9}
  \node[blk=notif]  (b1) at (90:\rad)  {\RL{\textbf{טריגר}}\\ \RL{התראה, רטט, נקודה אדומה}};
  \node[blk=notif]  (b2) at (45:\rad)  {\RL{\textbf{פעולה}}\\ \RL{פתיחת האפליקציה}};
  \node[blk=sun]    (b3) at (0:\rad)   {\RL{\textbf{תגמול משתנה}}\\ \RL{לייק? הודעה? סרטון?}};
  \node[blk=sun]    (b4) at (-45:\rad) {\RL{\textbf{השקעה}}\\ \RL{פוסט, תגובה, עוקבים}};
  \node[blk=screen] (b5) at (-90:\rad) {\RL{\textbf{איסוף נתונים}}\\ \RL{כל גלילה נרשמת}};
  \node[blk=screen] (b6) at (-135:\rad){\RL{\textbf{אלגוריתם}}\\ \RL{חיזוי מה יחזיק אותך}};
  \node[blk=night]  (b7) at (180:\rad) {\RL{\textbf{פרסום ממוקד}}\\ \RL{הקשב נמכר למפרסם}};
  \node[blk=night]  (b8) at (135:\rad) {\RL{\textbf{הכנסות}}\\ \RL{מימון פיצ'רים ממכרים}};

  % חיצים במעגל
  \foreach \a/\b in {b1/b2, b2/b3, b3/b4, b4/b5, b5/b6, b6/b7, b7/b8, b8/b1}
    \draw[arr] (\a) to[bend left=12] (\b);
  % חיצים מהמרכז
  \foreach \b in {b1,b3,b5,b7}
    \draw[-{Stealth}, dashed, draw=screen, line width=0.8pt] (core) -- (\b);

  % מספרי שלבים
  \foreach \b/\n/\c in {b1/1/notif, b2/2/notif, b3/3/sun, b4/4/sun,
                        b5/5/screen, b6/6/screen, b7/7/night, b8/8/night}
    \node[num=\c] at (\b.north east) {\n};
\end{tikzpicture}
\caption{מעגל כלכלת הקשב: שמונה שלבים שהופכים דקה של סקרנות לשעה של גלילה (עיבוד על בסיס מודל ה-\en{Hook} של \textcite{eyal2014} ו-\textcite{zuboff2019}).}
\label{fig:attention-loop}
\end{figure}


נלך במעגל, שלב אחר שלב:

\begin{enumerate}
  \item \textbf{טריגר}: הטלפון רוטט, נקודה אדומה מופיעה על האייקון. משהו קורא לכם.
  \item \textbf{פעולה}: אתם פותחים את האפליקציה. זה לוקח שנייה, ולא דורש מחשבה.
  \item \textbf{תגמול משתנה}: אולי לייק, אולי הודעה, אולי סרטון מצחיק. אי-הוודאות היא שמחזיקה אתכם.
  \item \textbf{השקעה}: אתם מגיבים, מעלים תמונה, עוקבים אחרי מישהו. כל השקעה כזאת מגדילה את הסיכוי שתחזרו.
  \item \textbf{איסוף נתונים}: כל גלילה, כל עצירה וכל הקלקה נרשמות.
  \item \textbf{אלגוריתם}: המערכת לומדת מה מחזיק דווקא אתכם, ומתאימה את התוכן הבא.
  \item \textbf{פרסום ממוקד}: הקשב שלכם, יחד עם מה שנלמד עליכם, נמכר למפרסמים.
  \item \textbf{הכנסות}: הכסף מושקע בחזרה בפיתוח פיצ'רים שמחזיקים אתכם עוד יותר, והמעגל מתחיל מחדש עם טריגר חדש.
\end{enumerate}

ארבעת השלבים הראשונים קורים אצלכם, בתודעה. ארבעת האחרונים קורים בשרתים, הרחק מהעיניים. ומה שהופך את המעגל לחזק כל כך הוא שכל סיבוב משפר את הסיבוב הבא. ככל שאתם גוללים יותר, האלגוריתם מכיר אתכם טוב יותר, וככל שהוא מכיר אתכם טוב יותר, קשה יותר להפסיק לגלול.

\begin{alertbox}[לא קונספירציה, מודל עסקי]
אין כאן חדר חשוך שבו מישהו זומם נגדכם. יש כאן מערכת שמכוונת להגדיל מדד אחד: זמן ומעורבות. כשמערכת חזקה מאוד ממקסמת מדד אחד, היא עלולה לפגוע בכל מה שאיננו נמדד, כמו שינה, ריכוז או שיחה בארוחת ערב.
\end{alertbox}

\section{קפיטליזם המעקב}

הסוציולוגית שושנה זובוף לוקחת את הטיעון צעד אחד קדימה. בספרה "עידן קפיטליזם המעקב" היא טוענת שהשיטה הזאת תובעת לעצמה את החוויה האנושית כ\textbf{חומר גלם חינמי} \parencite{zuboff2019}. החוויות שלנו מתורגמות לנתוני התנהגות. חלק מהנתונים משמש לשיפור השירות, אבל העודף, מה שזובוף מכנה \en{behavioral surplus}, הופך ל"מוצרי ניבוי": תחזיות על מה שנעשה, נקנה ונרגיש בהמשך.

כאן מסתתרת הקפיצה ההיסטורית. המהפכה התעשייתית הפכה את הטבע לחומר גלם: יערות, פחם, נפט. קפיטליזם המעקב, לפי זובוף, הופך את \textbf{הטבע האנושי} לחומר גלם. לא את מה שאנחנו מייצרים, אלא את מי שאנחנו.

\begin{thinkbox}
אם חברה יודעת מתי אתם עייפים, בודדים או משועממים, ויודעת גם מה תקנו ברגעים כאלה, מי בעצם מקבל את ההחלטה כשאתם לוחצים על "קנה עכשיו"?
\end{thinkbox}
